\documentclass[lualatex, paper={90mm, 160mm}]{jlreq}
\usepackage{tikz}
\usepackage{lltjext}
\usepackage{xcolor}
\usetikzlibrary{calc}

% 落ち着いたカラーパレット（くすみカラー）の設定
\definecolor{BgSand}{HTML}{F4F1EA}    % 背景：薄いサンドベージュ
\definecolor{MutedTeal}{HTML}{5D7B7C} % 曲線1：くすんだ青緑
\definecolor{DustyRose}{HTML}{B5838D} % 曲線2：くすんだローズ
\definecolor{SlateGray}{HTML}{6B705C} % 曲線3：スレートグレー

\begin{document}
\pagestyle{empty} 

\begin{tikzpicture}[remember picture, overlay]
  % 1. 背景全体を薄いベージュで塗りつぶす
  \fill[BgSand] (current page.south west) rectangle (current page.north east);

  % 2. 幾何学曲線（同心円と円弧のオーバーレイ）
  \begin{scope}[opacity=0.45, line width=0.8pt]
    % 左下からの波紋（くすんだ青緑）
    \foreach \r in {1, 3, ..., 15} {
      \draw[MutedTeal] ($(current page.south west) + (10mm, 15mm)$) circle (\r * 7mm);
    }
    % 右上からの波紋（くすんだローズ）
    \foreach \r in {2, 4, ..., 18} {
      \draw[DustyRose] ($(current page.north east) + (-15mm, -10mm)$) circle (\r * 8mm);
    }
    % 左上から右下へ流れる大きな円弧（スレートグレー）
    \foreach \r in {8, 11, 14} {
      \draw[SlateGray, line width=1.2pt] ($(current page.north west) + (0, 0)$) arc[start angle=180, end angle=300, radius=\r * 9mm];
    }
  \end{scope}

  % 3. 表紙の二重枠線（落ち着いたダークグレー）
  \draw[line width=1.2pt, color=black!75] ($(current page.south west) + (6mm, 6mm)$) rectangle ($(current page.north east) - (6mm, 6mm)$);
  \draw[line width=0.3pt, color=black!75] ($(current page.south west) + (7.5mm, 7.5mm)$) rectangle ($(current page.north east) - (7.5mm, 7.5mm)$);

  % 4. タイトル（縦書き） - 背景が明るいため白抜き背景は不要
  \node[anchor=north] at ($(current page.center) + (15mm, 50mm)$) {
    \begin{minipage}<t>{15em}
      \begin{center}
        {\fontsize{28pt}{40pt}\selectfont \textcolor{black!85}{走れメロス}\par}
      \end{center}
    \end{minipage}
  };

  % 5. 作者名（縦書き）
  \node[anchor=north] at ($(current page.center) + (-15mm, -10mm)$) {
    \begin{minipage}<t>{10em}
      \begin{center}
        {\fontsize{16pt}{24pt}\selectfont \textcolor{black!85}{太宰 治}\par}
      \end{center}
    \end{minipage}
  };
\end{tikzpicture}

\end{document}